\documentclass[12pt,a4paper]{article}

\usepackage[utf8]{inputenc}
\usepackage[margin=1in]{geometry}
\usepackage{graphicx}
\usepackage{setspace}
\usepackage{amsmath,amssymb}
\usepackage{booktabs}
\usepackage{hyperref}
\usepackage{cite}

\hypersetup{
    colorlinks=false,
    pdfborder={0 0 0}
}

\onehalfspacing

\begin{document}

% ==================== TITLE PAGE ====================
\thispagestyle{empty}
\begin{center}
    \vspace*{0.2cm}
    
    {\large \bfseries MULTI-SOURCE LINEAGE RETRIEVAL IN\\
    COMPOSITE MEDIA WITH DYNAMIC TOKEN\\
    ALIGNMENT AND IMPORTANCE GATING\par}
    
    \vspace{1.5cm}
    
    {\normalsize A Progress Report submitted for the degree of\par}
    \vspace{0.35cm}
    {\large \bfseries Master of Technology\par}
    \vspace{0.35cm}
    {\normalsize in\par}
    \vspace{0.35cm}
    {\large \bfseries Computer Science and Engineering\par}
    
    \vspace{1.4cm}
    
    {\normalsize Submitted by\par}
    \vspace{0.35cm}
    {\large \bfseries Vivek Negi\par}
    \vspace{0.2cm}
    {\normalsize \bfseries Roll Number: T25CS012\par}
    
    \vspace{1.2cm}
    
    {\normalsize Under the Supervision of\par}
    \vspace{0.35cm}
    {\large \bfseries Dr. Akhilendra Pratap Singh\par}
    \vspace{0.2cm}
    {\normalsize Assistant Professor\par}
    
    \vspace{1.0cm}
    
    % Automatically renders the NITM logo if uploaded (as nitm_logo.png, logo.jpg, or logo.png),
    % otherwise leaves a clean 3.2cm space so Overleaf compiles without errors!
    \IfFileExists{nitm_logo.png}{%
        \includegraphics[width=3.2cm]{nitm_logo.png}\par
    }{%
        \IfFileExists{logo.jpg}{%
            \includegraphics[width=3.2cm]{logo.jpg}\par
        }{%
            \IfFileExists{logo.png}{%
                \includegraphics[width=3.2cm]{logo.png}\par
            }{%
                \vspace{3.2cm}
            }%
        }%
    }
    
    \vfill
    
    {\normalsize \bfseries Department of Computer Science and Engineering\par}
    {\normalsize \bfseries National Institute of Technology Meghalaya\par}
    {\normalsize \bfseries Saitsohpen Sohra--793108\par}
    \vspace{0.4cm}
    {\normalsize November, 2026\par}
\end{center}

% ==================== DECLARATION ====================
\newpage
\pagenumbering{roman}
\setcounter{page}{1}

\begin{center}
    {\Large \bfseries Declaration\par}
\end{center}

\vspace{1.2cm}

\noindent I hereby declare that the End-Semester Progress Report entitled \textbf{``Multi-Source Lineage Retrieval in Composite Media with Dynamic Token Alignment and Importance Gating''} submitted towards the fulfilment of the Master of Technology dissertation review in the Department of Computer Science and Engineering at the National Institute of Technology Meghalaya is an authentic record of research work carried out by me under the supervision of \textbf{Dr. Akhilendra Pratap Singh} and has not been submitted elsewhere for any academic degree or qualification.

\vspace{1.0cm}

\noindent I will be solely responsible if any kind of misrepresentation is found in the report.

\vspace{2.5cm}

\begin{flushright}
    \textbf{VIVEK NEGI} \\
    Roll No: T25CS012 \\
    Department of CSE, NIT Meghalaya
\end{flushright}

% ==================== TABLE OF CONTENTS ====================
\newpage
\tableofcontents
\newpage

\pagenumbering{arabic}
\setcounter{page}{1}

% ==================== SECTION 1 ====================
\section{Introduction}
Rapid growth of multimedia sharing platforms such as Instagram, YouTube, X (Twitter), and decentralized digital networks has significantly increased the creation, remixing, and redistribution of digital content. Users continuously upload and modify videos, audio clips, images, memes, and edited photographic media across online platforms. Although this digital expansion improves communication, creative expression, and accessibility, it also increases digital copyright infringement, automated insurance fraud, and criminal evidence tampering challenges \cite{sscd, trufor}.

Modern user-generated multimedia manipulation frequently involves multi-source compositing (image splicing). Formally, a composite query image $Q$ is created by combining two distinct visual sources:
\begin{equation}
Q = M \odot I_A + (1 - M) \odot I_B
\end{equation}
where $M \in [0, 1]^{H \times W}$ is the blending mask and $\odot$ represents pixel-wise multiplication. In such a composite image:
\begin{enumerate}
    \item \textbf{Parent A (Background Scene):} Serves as the main background photo, typically occupying 70\% to 85\% of the total image area.
    \item \textbf{Parent B (Cut-and-Pasted Donor Object):} Represents a smaller object (such as a person, handgun, license plate, or watch) cut out from a second photo and pasted onto the background scene (occupying 10\% to 30\% of the image area).
\end{enumerate}

Existing reverse-image search systems, such as Google Lens, TinEye, or traditional perceptual hashing (pHash, dHash), assume that an edited photo comes from only \textit{one} original photo. When given a spliced composite photo, these systems easily find the big background (Parent A) but completely miss the smaller pasted object (Parent B).

The core goal of \textbf{Multi-Source Lineage Retrieval} is to build an automated forensic system that takes a composite photo $Q$, searches a database gallery $\mathcal{G} = \{I_1, I_2, \dots, I_N\}$ of original photos, and retrieves \textit{both} Parent A and Parent B together while verifying on a blockchain which photo was created first.

% ==================== SECTION 2 ====================
\section{Literature Review}
This section reviews prior works in self-supervised copy detection, image splicing localization, patch-level feature extraction, generative benchmarks, and blockchain media provenance.

\subsection{Deep Feature Descriptors for Copy Detection}
Pizzi et al. (CVPR 2022) \cite{sscd} introduced the Self-Supervised Descriptor for Copy Detection (SSCD), developed at Meta AI. Built on a ResNet-50 convolutional neural network backbone with Generalized Mean (GeM) pooling, SSCD converts an entire image into a 512-dimensional feature vector. While SSCD is very strong at finding cropped, recolored, or compressed copies of a single image, compressing the whole image into one global vector causes small pasted donor objects to be washed out by the large background.

\subsection{Forensic Splicing Localization}
Guillaro et al. (CVPR 2023) \cite{trufor} proposed TruFor, a forensic network that detects and highlights spliced regions in an image by analyzing camera sensor noise traces \cite{noiseprint} and compression artifacts. However, TruFor only outputs a heatmap showing \textit{where} an edit occurred; it cannot search a database to tell us \textit{which original photo} the pasted object came from.

\subsection{Patch-Level Sensor Artifact Analysis}
Chen, Yao, and Niu (2024) \cite{singlepatch} showed in \textit{``A Single Simple Patch is All You Need for AI-generated Image Detection''} that examining small local image patches preserves fine forensic details that whole-image pooling loses. This supports our approach of extracting features from local image patches rather than only the whole image.

\subsection{Generative Benchmarking on Cropping vs. Resizing}
Li et al. (NeurIPS 2025) \cite{aigibench} introduced AIGIBench and proved that shrinking/resizing a full image destroys fine-grained textures, whereas cropping local windows preserves authentic visual features. This directly supports our multi-scale sliding window cropping design.

\subsection{Cryptographic Provenance and Metadata Stripping}
Wazid et al. (2024) \cite{wazid} and the C2PA standard \cite{c2pa} attach cryptographic signatures inside image EXIF metadata headers. However, social media apps (WhatsApp, Instagram, X) automatically delete EXIF metadata when an image is uploaded. By combining visual feature matching with an external SHA-256 blockchain ledger, our system recovers provenance even when file metadata is stripped. Table~\ref{tab:literature} summarizes the related literature.

\begin{table}[h]
\centering
\small
\begin{tabular}{p{3.0cm} p{3.6cm} p{3.5cm} p{3.8cm}}
\toprule
\textbf{Study} & \textbf{Contribution} & \textbf{Technique} & \textbf{Key Limitation} \\
\midrule
Meta SSCD \cite{sscd} & Copy detection SOTA & ResNet-50 + GeM, 512-d & Global vector misses donors \\
TruFor \cite{trufor} & Splicing heatmap & Noiseprint++ transformer & No database search \\
Single Patch \cite{singlepatch} & Local patch analysis & High-pass patch filters & Fake/real classifier only \\
AIGIBench \cite{aigibench} & Benchmark study & Crop vs. resize analysis & Evaluation dataset only \\
C2PA / Wazid \cite{wazid} & Metadata signatures & EXIF header manifests & Stripped by social apps \\
\textbf{ProvenanceAI} & \textbf{Dual-Parent + Chain} & \textbf{Dynamic Crops + SHA-256} & \textbf{Finds both A \& B (88.89\%)} \\
\bottomrule
\end{tabular}
\caption{Comparative Summary of Related Works and ProvenanceAI}
\label{tab:literature}
\end{table}

% ==================== SECTION 3 ====================
\section{Research Gap and Motivation}

\subsection{Research Gap: The Background Dominance Problem}
Standard neural networks convert a full image into one single feature vector $\mathbf{v}_{\text{global}}$. In a spliced photo, because the background (Parent A) covers 70\% to 85\% of the pixels, the global vector is dominated by the background:
\begin{equation}
\cos(\mathbf{v}_Q, \mathbf{v}_{\text{Parent A}}) \approx 0.96 \quad \text{vs.} \quad \cos(\mathbf{v}_Q, \mathbf{v}_{\text{Parent B}}) \le 0.42
\end{equation}
As a result, standard search retrieves Parent A (100\% recall) but fails on Parent B (only 44.44\% recall). Furthermore, cutting the image into a fixed $2\times2$ or $3\times3$ grid slices objects standing in the middle of the photo into separate pieces, ruining recognition.

\subsection{Motivation}
Our motivation is two-fold: (1) accurately recover both the background photo and the small pasted object without expensive retraining of large vision models, and (2) record every original photo and audit result on a tamper-proof SHA-256 blockchain so that no one can backdate a fake composite photo to claim it was the original.

% ==================== SECTION 4 ====================
\section{Proposed Framework and Implementation Stack}
Our system, \textbf{ProvenanceAI}, consists of four simple, modular stages backed by a local Python software stack.

\subsection{Stage 1: Multi-Scale Overlapping Window Slicing}
Instead of a rigid grid that cuts centered objects in half, we slide overlapping crop windows across the query image at three window sizes $S = \{25\%, 35\%, 50\%\}$ of the image dimensions with a \textbf{50\% overlap step}:
\begin{equation}
\text{stride}_x = \max\left(20, \left\lfloor \frac{w_s}{2} \right\rfloor\right), \quad \text{stride}_y = \max\left(20, \left\lfloor \frac{h_s}{2} \right\rfloor\right)
\end{equation}
Because windows overlap by 50\%, at least one window is guaranteed to capture the pasted object (Parent B) cleanly in one piece.

\subsection{Stage 2: Lightweight MLP Importance Gate}
Many cropped windows only contain empty sky, plain walls, or flat road, which can cause false matches. We pass each patch vector $\mathbf{p}_i \in \mathbb{R}^{128}$ through a lightweight 2-layer Multi-Layer Perceptron (MLP) neural gate:
\begin{equation}
w_i = \sigma\left(\mathbf{W}_2 \cdot \text{ReLU}(\mathbf{W}_1 \mathbf{p}_i + \mathbf{b}_1) + \mathbf{b}_2\right)
\end{equation}
where $\sigma(\cdot)$ outputs an importance weight $w_i \in [0, 1]$. Plain background patches receive low weights ($w_i \to 0$), while detailed foreground objects (faces, cars, weapons, watches) receive high weights ($w_i \to 1$). This gate has fewer than 70,000 parameters and runs in milliseconds.

\subsection{Stage 3: Dual-Channel Matching Score ($\alpha = 0.25$)}
For every image $G_k$ in our database, we combine the whole-image match score and the best weighted patch match score:
\begin{equation}
\mathcal{S}_{\text{lineage}}(Q, G_k) = \alpha \cdot S_{\text{global}}(Q, G_k) + (1 - \alpha) \cdot \max_{i} \left( w_i \cdot \cos(\mathbf{p}_i^Q, \mathbf{f}_{\text{global}}^{G_k}) \right)
\end{equation}
We set $\alpha = 0.25$ (giving 25\% weight to the global view and 75\% weight to the best local patch). Parent A is identified by the highest global similarity, and Parent B is identified by the highest gated patch similarity among the remaining database images.

\subsection{Stage 4: Custom SHA-256 Proof-of-Provenance Blockchain}
Every original gallery photo and verified composite is recorded on our custom Python blockchain ledger (\texttt{blockchain.py} and \texttt{blockchain\_ledger.json}), which currently holds \textbf{129 mined blocks}. Each block stores the block index, creation timestamp, image SHA-256 file hash, feature embedding hash, parent lineage pointers, Proof-of-Work nonce, and previous block hash. The ledger verifies that both Parent A and Parent B existed on-chain \textit{before} the composite query $Q$:
\begin{equation}
\text{Verified}(Q) \iff \mathcal{T}(\text{Parent A}) < \mathcal{T}(Q) \;\land\; \mathcal{T}(\text{Parent B}) < \mathcal{T}(Q)
\end{equation}

\subsection{Software Architecture and Tools Used}
ProvenanceAI runs 100\% locally on the machine so sensitive forensic evidence never has to be uploaded to third-party cloud APIs. Our built-in Forensic Splicing Studio uses OpenCV and Pillow for clean alpha cutouts, Reinhard color matching, and a \textbf{Dual-Layer Contact Shadow Engine} that casts realistic surface shadows under flat objects (handguns, watches) and standing people. Table~\ref{tab:stack} lists the tools used.

\begin{table}[h]
\centering
\small
\begin{tabular}{p{3.2cm} p{4.2cm} p{6.8cm}}
\toprule
\textbf{Subsystem Layer} & \textbf{Language / Tools} & \textbf{Role in ProvenanceAI} \\
\midrule
Core Language & Python 3.9+, NumPy & Fast matrix math, sliding windows, cosine similarity \\
Deep Vision Model & PyTorch, ResNet-50 / SSCD & Frozen 512-d feature extractor \& 2-layer MLP gate \\
Training Hardware & NVIDIA Tesla T4 GPU & Trains the $<70\text{k}$ param MLP gate in $<1$ minute \\
Image Processing & Pillow (PIL), OpenCV & Object cutouts, color matching, dual-layer shadows \\
Blockchain Ledger & Python \texttt{hashlib} (SHA-256) & Proof-of-Work block mining \& 129-block JSON ledger \\
Web Dashboard & Streamlit Framework & 5-tab interactive forensic UI on \texttt{localhost:8502} \\
Live Public Link & SSH SSL Reverse Tunnel & Remote access at \url{https://uuunw-115-245-187-126.run.pinggy-free.link} \\
\bottomrule
\end{tabular}
\caption{Software Engineering Stack and Tools Used in ProvenanceAI}
\label{tab:stack}
\end{table}

% ==================== SECTION 5 ====================
\section{Failure Mode Analysis: AI Diffusion Edits and FFT Gating}

\subsection{Where Standard Retrieval Fails: Generative AI Inpainting}
We also tested what happens when an image is edited using a generative AI diffusion model (such as Meta AI Emu or Stable Diffusion \cite{ldm}) rather than simple cut-and-paste. In our test, a user gave two photos to an AI model and prompted it to blend the sun from Photo 2 into the sky of Photo 1 (\texttt{vk.png}).

Our system found Parent A (\texttt{vk.png}, Block 62) with 89.4\% confidence, but initially matched the AI-generated sun to a yellow \texttt{Studio Headphones} photo in our database (81.86\% match) simply because both were bright yellow circles. \textbf{Why this happens:} Generative diffusion models do not copy-paste actual pixels from Photo 2; they redraw brand-new pixels from random Gaussian noise. Therefore, no real camera pixels from Photo 2 exist inside the output image.

\subsection{Our Fix: 2D FFT Frequency Filter}
To stop the system from guessing a false Parent B on AI-redrawn regions, we added a 2D Fast Fourier Transform (FFT) frequency check on each candidate patch:
\begin{equation}
F(u, v) = \sum_{x=0}^{M-1} \sum_{y=0}^{N-1} I(x, y) e^{-j 2\pi \left(\frac{ux}{M} + \frac{vy}{N}\right)}
\end{equation}
Real camera photos contain sharp high-frequency sensor grain, whereas AI diffusion smoothing wipes out high-frequency energy. When our FFT filter sees abnormally low high-frequency energy ($\rho_{\text{HF}} < 0.88$ or $\sigma_{\text{HF}} < 14.5$), it automatically blocks the false match and flags the region as: \textit{``Synthetic AI Inpainting Detected (No Physical Donor)''}, while still reporting Parent A accurately.

% ==================== SECTION 6 ====================
\section{Experimental Evaluation and Dataset Results}

\subsection{Dataset Size (MS-COCO 2017) and Hardware Setup}
We trained and evaluated ProvenanceAI using real-world photographs across two dataset tiers:
\begin{enumerate}
    \item \textbf{MS-COCO 2017 Benchmark Dataset:} Built from official \textbf{MS-COCO 2017} natural photographs (people, vehicles, animals, indoor/outdoor scenes) \cite{coco}. The search database consists of \textbf{1,000 unindexed gallery photos} and \textbf{300 realistic composite queries} created with Gaussian edge blending (donor size between 15\% and 35\% of image area). The 300 composites are split into:
    \begin{itemize}
        \item \textbf{Training Set (70\%):} 210 composite queries used to train our 2-layer MLP Importance Gate in under 1 minute on an \textbf{NVIDIA Tesla T4 GPU}.
        \item \textbf{Validation Set (15\%):} 45 composite queries used to tune $\alpha = 0.25$.
        \item \textbf{Held-Out Test Set (15\%):} 45 composite queries used for final benchmark evaluation.
    \end{itemize}
    \item \textbf{Real-World Forensic Gallery (Live App):} Over \textbf{40+ high-resolution forensic photographs} (urban crosswalks, vehicles, Indian HSRP license plates, crime scene desks, university libraries, studio portraits) notarized across \textbf{129 blockchain blocks} and tested across 6 realistic forensic fraud scenarios.
\end{enumerate}

\subsection{Benchmark Accuracy and Ablation Study}
Table~\ref{tab:results} compares our full pipeline against standard baseline methods.

\begin{table}[h]
\centering
\small
\begin{tabular}{l c c c c}
\toprule
\textbf{Method / Pipeline Strategy} & \textbf{Parent A Recall} & \textbf{Parent B Recall} & \textbf{Full Lineage@5} & \textbf{Latency} \\
\midrule
Global SSCD Baseline \cite{sscd} & 100.00\% & 44.44\% & 44.44\% & 190 ms \\
Fixed $2\times2$ Grid Split & 100.00\% & 66.67\% & 66.67\% & 340 ms \\
Fixed $3\times3$ Grid Split & 100.00\% & 66.67\% & 66.67\% & 410 ms \\
Dynamic Overlapping Crops (No Gate) & 100.00\% & 77.78\% & 77.78\% & 365 ms \\
\textbf{ProvenanceAI (Dynamic + MLP Gate)} & \textbf{100.00\%} & \textbf{88.89\%} & \textbf{88.89\%} & \textbf{380 ms} \\
\bottomrule
\end{tabular}
\caption{Full-Lineage Retrieval Recall@5 Benchmark Comparison}
\label{tab:results}
\end{table}

Overlapping windows boosted Parent B recovery from 44.44\% to 77.78\% by keeping centered objects whole, and adding the 2-layer MLP Importance Gate pushed final accuracy to \textbf{88.89\% Full Lineage Recall@5} (+44.45\% gain over global SSCD) in just 380 ms per query.

\subsection{Validation Across Six Real-World Forensic Case Studies}
\begin{enumerate}
    \item \textbf{Case 1 (Mountain Alibi Fabrication):} Person cutout pasted into mountain overlook $\rightarrow$ Retrieved Parent A (99.99\%) and Parent B (92.65\%), verified in Block 125.
    \item \textbf{Case 2 (Vehicle License Plate Swap):} Indian HSRP plate (\texttt{DL 01 AB 4921}) spliced onto sports car $\rightarrow$ 100\% retrieval of both original car and donor plate.
    \item \textbf{Case 3 (Crime Scene Evidence Planting):} 9mm handgun pasted onto empty wooden desk $\rightarrow$ Retrieved Parent A (99.98\%) and Parent B (96.88\%).
    \item \textbf{Case 4 (Urban Crosswalk Hit-and-Run):} Pedestrian pasted into city crosswalk $\rightarrow$ Retrieved Parent A (99.99\%) and Parent B (93.45\%).
    \item \textbf{Case 5 (University Library Contraband):} Leather backpack spliced into library hall (\texttt{real\_modern\_library\_hall.jpg}) $\rightarrow$ 99.97\% Parent A, 92.40\% Parent B (Block 127).
    \item \textbf{Case 6 (Workspace Watch Audit):} Chronograph watch spliced onto office desk (\texttt{real\_wooden\_work\_desk.jpg}) $\rightarrow$ 99.98\% Parent A, 98.35\% Parent B (Block 128).
\end{enumerate}

% ==================== SECTION 7 ====================
\section{Objectives Accomplished}
\begin{itemize}
    \item Solved the Background Dominance Problem in composite photos without retraining heavy neural networks, doubling dual-parent retrieval accuracy from 44.44\% to 88.89\%.
    \item Built the Multi-Scale Dynamic Overlapping Slicer and 2-layer MLP Importance Gate ($<70\text{k}$ parameters) in PyTorch.
    \item Created an interactive Forensic Splicing Studio with realistic dual-layer contact shadows and tested it on 40+ real photographs across 6 crime/fraud scenarios.
    \item Added a 2D FFT frequency filter to detect AI diffusion edits and prevent false donor matches.
    \item Implemented a 129-block SHA-256 Proof-of-Provenance Blockchain Ledger and live Streamlit web app accessible locally and via remote tunnel.
\end{itemize}

% ==================== SECTION 8 ====================
\section{Conclusion and Future Scope}
This report presented \textbf{ProvenanceAI}, a practical system that finds both the background source and the pasted donor object in manipulated composite photos while proving creation timestamps on a SHA-256 blockchain. Combining overlapping window crops, a lightweight MLP importance gate ($\alpha=0.25$), and a 2D FFT frequency filter achieves 88.89\% Full-Lineage Recall@5 in 380 ms. The complete system is deployed at \url{http://localhost:8502} and externally accessible via \url{https://uuunw-115-245-187-126.run.pinggy-free.link}.

For future scope, simple modular extensions include: (1) running our sliding-window matcher on extracted keyframes from short video clips, (2) quantizing the MLP gate to 8-bit integers (INT8) for mobile deployment, and (3) adding zoomable lineage tree graphs in the web UI.

% ==================== REFERENCES ====================
\newpage
\begin{thebibliography}{99}
\bibitem{sscd} E. Pizzi, S. Durasamy, S. Couprie, M. Douze, and A. Gordo, ``A self-supervised descriptor for image copy detection,'' in \textit{Proceedings of the IEEE/CVF Conference on Computer Vision and Pattern Recognition (CVPR)}, 2022, pp. 14532--14542.
\bibitem{trufor} F. Guillaro, D. Cozzolino, A. Sud, N. Dufour, and L. Verdoliva, ``TruFor: Leveraging all-round clues for trustworthy image forgery detection and localization,'' in \textit{Proceedings of the IEEE/CVF Conference on Computer Vision and Pattern Recognition (CVPR)}, 2023, pp. 9601--9610.
\bibitem{singlepatch} J. Chen, J. Yao, and L. Niu, ``A single simple patch is all you need for AI-generated image detection,'' \textit{arXiv preprint arXiv:2402.01123}, 2024.
\bibitem{aigibench} Z. Li, J. Yan, Z. He, K. Zeng, et al., ``Is artificial intelligence generated image detection a solved problem? (AIGIBench),'' in \textit{Advances in Neural Information Processing Systems (NeurIPS)}, 2025.
\bibitem{wazid} M. Wazid, A. K. Das, and Y. Park, ``Blockchain-based secure multimedia copyright protection framework,'' \textit{IEEE Open Journal of the Computer Society}, vol. 5, pp. 248--267, 2024.
\bibitem{c2pa} Coalition for Content Provenance and Authenticity (C2PA), ``Content Credentials: Technical specifications for digital media provenance,'' \textit{C2PA Standards Whitepaper}, 2024.
\bibitem{noiseprint} D. Cozzolino and L. Verdoliva, ``Noiseprint: A CNN-based camera model fingerprint,'' \textit{IEEE Transactions on Information Forensics and Security}, vol. 15, pp. 377--389, 2020.
\bibitem{moreira} P. Moreira, A. Bharati, et al., ``Image provenance analysis at scale,'' \textit{IEEE Transactions on Information Forensics and Security}, vol. 13, no. 12, pp. 2857--2871, 2018.
\bibitem{dias} Z. Dias, A. Rocha, and S. Goldenstein, ``Image phylogeny by minimal spanning trees,'' \textit{IEEE Transactions on Pattern Analysis and Machine Intelligence}, vol. 34, no. 12, pp. 2362--2375, 2012.
\bibitem{ldm} R. Rombach, A. Blattmann, D. Lorenz, P. Esser, and B. Ommer, ``High-resolution image synthesis with latent diffusion models,'' in \textit{Proceedings of the IEEE/CVF Conference on Computer Vision and Pattern Recognition (CVPR)}, 2022, pp. 10684--10695.
\bibitem{coco} T.-Y. Lin, M. Maire, S. Belongie, J. Hays, et al., ``Microsoft COCO: Common objects in context,'' in \textit{Proceedings of the European Conference on Computer Vision (ECCV)}, 2014, pp. 740--755.
\end{thebibliography}

\end{document}
